\documentclass[12pt]{article}

\usepackage{sbc-template}
\usepackage{graphicx,url}
\usepackage[brazilian]{babel}
\usepackage[utf8]{inputenc}
\usepackage[T1]{fontenc}
\usepackage{booktabs}
\usepackage{array}

\sloppy

\title{Aplicação do modelo DMC ao mini jogo Travessia do Rio: adaptação plugada de uma atividade desplugada}

\author{Filipe Sales\inst{1}}

\address{Plataforma TEA -- Ensino de programação e Pensamento Computacional\\
  para estudantes com Transtorno do Espectro Autista
}

\begin{document}

\maketitle

\begin{abstract}
  This paper records how the DMC model (Dynamics, Mechanics and Components) was
  applied to the Travessia do Rio mini game of the Plataforma TEA. The game is a
  digital (plugged) adaptation of Task 1, an unplugged card game used with
  Problem-Based Learning (PBL) to teach the idea of algorithm. The application is
  deliberately basic: most elements were already present in the game and were
  only named and mapped, and one new component, a delivery progress indicator,
  was added. Points, rankings and timers were left out on purpose, to respect the
  sensory and pedagogical rules of the project. No evaluation with students has
  been carried out yet.
\end{abstract}

\begin{resumo}
  Este trabalho registra como o modelo DMC (Dinâmicas, Mecânicas e Componentes)
  foi aplicado ao mini jogo Travessia do Rio da Plataforma TEA. O jogo é uma
  adaptação digital (plugada) da Tarefa 1, um jogo de cartas desplugado usado com
  Aprendizagem Baseada em Problemas (ABP) para trabalhar a ideia de algoritmo. A
  aplicação é deliberadamente básica: a maior parte dos elementos já existia no
  jogo e foi apenas nomeada e mapeada, e um componente novo, o indicador de
  progresso de entregas, foi acrescentado. Pontos, ranqueamento e cronômetro
  ficaram de fora de propósito, para respeitar as regras sensoriais e
  pedagógicas do projeto. Ainda não houve avaliação com estudantes.
\end{resumo}

\section{Introdução}

A Plataforma TEA é uma plataforma web gamificada para ensinar lógica de
programação e Pensamento Computacional a estudantes com Transtorno do Espectro
Autista (TEA). O mini jogo Travessia do Rio é um dos jogos autônomos da
plataforma, voltado à ideia de algoritmo como passo a passo: repetição, decisão
e verificação.

O jogo não nasceu do zero. Ele é a versão plugada da Tarefa 1 do conjunto de
atividades do projeto, uma atividade desplugada \cite{bell2009} conduzida com
Aprendizagem Baseada em Problemas (ABP) \cite{barrows1986,hmelo2004}, descrita no
plano de aula \texttt{travessia-rio.tex} (diretório \texttt{docs/desplugadas}).
Este documento registra a aplicação do modelo DMC \cite{werbach2012} a essa
adaptação, no diretório \texttt{docs/plugadas}.

\section{O modelo DMC}

O modelo DMC organiza a gamificação em três níveis, do mais abstrato ao mais
concreto \cite{werbach2012}:

\begin{itemize}
  \item \textbf{Dinâmicas}: aspectos gerais do sistema, que não aparecem
        diretamente na tela, como restrições, progressão e narrativa;
  \item \textbf{Mecânicas}: os processos que movem a ação e o engajamento, como
        desafios, recursos, \textit{feedback} e turnos;
  \item \textbf{Componentes}: as instâncias concretas das dinâmicas e mecânicas,
        como níveis, avatares e indicadores.
\end{itemize}

Gamificação, aqui, é o uso de elementos de design de jogos em contextos que não
são jogos \cite{deterding2011}. No projeto, esses elementos passam antes por um
filtro sensorial: o que não tem respaldo em evidência nasce desligado ou fica de
fora.

\section{Da atividade desplugada ao jogo plugado}

A Tarefa 1 usa tabuleiro, peças e fluxograma em papel. Cinco personagens com
pesos diferentes (Lobo, Ovelha, Alface, Cachorro e Galinha) atravessam um rio em
um barco de peso máximo 3, conduzido pelo Fazendeiro, sob quatro regras de
conflito (R1 a R4) e com um tanque de combustível limitado. A rodada 1 é um
aquecimento sem o Lobo, com 7 gotas, e a rodada 2 é o jogo completo, com 13
gotas. As melhores soluções têm 5 e 11 travessias, respectivamente.

A Tabela~\ref{tab:adaptacao} resume o que foi mantido e o que mudou na
adaptação.

\begin{table}[ht]
\centering
\small
\caption{O que a versão plugada herda da Tarefa 1 e o que muda.}
\label{tab:adaptacao}
\begin{tabular}{@{}>{\raggedright\arraybackslash}p{3.0cm}>{\raggedright\arraybackslash}p{3.4cm}>{\raggedright\arraybackslash}p{4.6cm}@{}}
\toprule
\textbf{Elemento} & \textbf{Desplugada} & \textbf{Plugada} \\
\midrule
Regras e rodadas & R1 a R4, rodadas 1 e 2 & Mantidas, com os mesmos números \\
Verificação & O fiscal confere a margem sem o barco & O jogo confere sozinho e desfaz a travessia inválida \\
Combustível & Gotas riscadas à mão & Tanque na tela; a gota gasta continua gasta \\
Dicas & Cartas entregues pelo docente, com desconto na nota & Cartas virtuais; o desconto só é registrado em evento \\
Papéis e Ficha & Operador, fiscal, anotador e Ficha & Sem papéis; o registro vira eventos estruturados \\
Fluxograma & Produto do grupo & Fora do jogo; segue no material desplugado \\
\bottomrule
\end{tabular}
\end{table}

\section{Aplicação do DMC}

A aplicação do DMC foi mantida simples. Primeiro, verificou-se o que o jogo já
tinha: ele não declarava o modelo, mas várias mecânicas já existiam, herdadas da
Tarefa 1. O trabalho consistiu em mapear esses elementos nos três níveis e
acrescentar um componente que faltava.

\subsection{Dinâmicas}

\begin{itemize}
  \item \textbf{Restrição:} o barco tem peso máximo e o combustível é limitado.
        É a restrição que transforma o enigma em problema de planejamento.
  \item \textbf{Progressão:} a passagem da rodada 1 para a rodada 2 aumenta a
        complexidade (entra o Lobo, o tanque cresce) e é medida só em relação ao
        que o próprio estudante já fez.
  \item \textbf{Narrativa:} a fazenda e o Fazendeiro dão contexto sem texto longo.
        Os personagens são desenhos; o nome vai para o leitor de tela.
\end{itemize}

\subsection{Mecânicas}

\begin{itemize}
  \item \textbf{Desafio:} levar todos à margem direita sem violar R1 a R4.
  \item \textbf{Recurso:} as gotas de combustível, gastas a cada travessia,
        inclusive as desfeitas.
  \item \textbf{\textit{Feedback}:} o conflito aparece na margem, com uma frase
        descritiva e reversível (``Travessia desfeita. Quer tentar outra
        carga?''). Não há ``errado'', ``falhou'' nem sinal vermelho grande.
  \item \textbf{Ajuda:} cartas de dica viradas sob demanda, que ficam abertas
        para serem relidas.
  \item \textbf{Turnos e tentativa:} cada travessia é uma jogada; é possível
        recomeçar a rodada quando se quiser.
\end{itemize}

\subsection{Componentes}

\begin{itemize}
  \item \textbf{Níveis:} as duas rodadas.
  \item \textbf{Avatares:} os cinco personagens, com o peso em um número no canto.
  \item \textbf{Tanque de combustível:} indicador de gotas restantes.
  \item \textbf{Cartas de dica:} três na rodada 1 e cinco na rodada 2.
  \item \textbf{Indicador de progresso de entregas (novo):} mostra quantos
        personagens já estão na margem direita (por exemplo, ``2 de 5''). Sobe
        e desce conforme a carga atravessa em um sentido ou no outro, e uma
        travessia desfeita não conta como entrega.
\end{itemize}

A Tabela~\ref{tab:dmc} reúne o mapa completo.

\begin{table}[ht]
\centering
\small
\caption{DMC aplicado à Travessia do Rio.}
\label{tab:dmc}
\begin{tabular}{@{}>{\raggedright\arraybackslash}p{2.4cm}>{\raggedright\arraybackslash}p{2.8cm}>{\raggedright\arraybackslash}p{4.0cm}@{}}
\toprule
\textbf{Dinâmica} & \textbf{Mecânica} & \textbf{Componente} \\
\midrule
Restrição & Recurso limitado & Tanque de combustível; peso do barco \\
Restrição & Desafio com regras & Regras de conflito R1 a R4 \\
Progressão & Níveis & Rodadas 1 e 2 \\
Progressão & \textit{Feedback} de avanço & Indicador de entregas (novo) \\
Narrativa & Personagens & Avatares na fazenda \\
Apoio & Ajuda sob demanda & Cartas de dica \\
Apoio & \textit{Feedback} não punitivo & Mensagem de conflito e desfazer \\
\bottomrule
\end{tabular}
\end{table}

\subsection{O que ficou de fora de propósito}

O projeto tem regras que limitam a gamificação, em linha com a evidência sobre
TEA que orienta a plataforma. Por isso, não foram aplicados:

\begin{itemize}
  \item \textbf{pontos, \textit{badges} e ranqueamento}: não há comparação entre
        estudantes, e a progressão é sempre relativa ao próprio histórico. O
        jogo continua sem pontuação, como a atividade desplugada (a ``nota'' era
        só do docente);
  \item \textbf{cronômetro e contagem regressiva}: a pressão de tempo não existe
        no jogo;
  \item \textbf{efeitos extras}: nenhuma animação ou som novo foi acrescentado;
        o indicador de progresso é estático.
\end{itemize}

\section{Registro de dados}

O jogo já registra eventos estruturados quando há aluno logado (início de
rodada, travessias, carga, dicas, conclusão e abandono). O indicador de entregas
é derivado do estado do jogo e não gera evento próprio. O desconto de cada dica,
que na atividade desplugada reduzia a nota, é gravado no evento de pedido de
dica para o docente cruzar depois, sem aparecer para o estudante.

\section{Considerações finais}

O modelo DMC foi aplicado de forma básica à Travessia do Rio: os elementos
existentes foram nomeados e mapeados, e o único componente novo foi o indicador
de progresso de entregas. A adaptação preserva as regras, os números e a lógica
da Tarefa 1 e acrescenta a verificação automática e o \textit{feedback} não
punitivo.

Como limitações, a aplicação não foi testada com estudantes com TEA, e os
efeitos da gamificação sobre engajamento e aprendizagem não foram medidos. Como
trabalho futuro, propõem-se um teste de uso em turma e a avaliação do
indicador de progresso com os registros de eventos já coletados.

\bibliographystyle{sbc}
\begin{thebibliography}{9}

\bibitem{barrows1986}
Barrows, H. S. (1986). A taxonomy of problem-based learning methods.
\textit{Medical Education}, 20(6):481--486.

\bibitem{bell2009}
Bell, T., Alexander, J., Freeman, I., and Grimley, M. (2009). Computer science
unplugged: school students doing real computing without computers.
\textit{The New Zealand Journal of Applied Computing and Information
Technology}, 13(1):20--29.

\bibitem{deterding2011}
Deterding, S., Dixon, D., Khaled, R., and Nacke, L. (2011). From game design
elements to gamefulness: defining ``gamification''. In \textit{Proceedings of the
15th International Academic MindTrek Conference}, pages 9--15. ACM.

\bibitem{hmelo2004}
Hmelo-Silver, C. E. (2004). Problem-based learning: what and how do students
learn? \textit{Educational Psychology Review}, 16(3):235--266.

\bibitem{werbach2012}
Werbach, K. and Hunter, D. (2012). \textit{For the Win: How Game Thinking Can
Revolutionize Your Business}. Wharton Digital Press, Philadelphia.

\end{thebibliography}

\end{document}
